\vfill%
\pagebreak
\section{Conditional branching}%
\label{sec:flow:if_else}

A conditional evaluates one of two branches, depending on a condition of type
\token{bool}. Its variant \token{if let} depends instead on whether a value fits
a pattern (Section~\ref{sec:flow:cond_pattern}), and \token{cte if} is resolved
at compile time (Section~\ref{sec:flow:cte_if}).

\subsection{Grammar}

\begin{lstlisting}[style=grammarVerb]
if_expression := 'cte'? 'if' condition body else?
               | 'if' 'let' pattern '=' expression guard? body else?
else          := 'else' body
               | 'else' if_expression
body          := block
               | expression ';'?
guard         := '&&' expression
\end{lstlisting}

A \token{body} that is not a block is a single expression, optionally followed
by a semicolon: \token{if c println("a"); else println("b");} is valid. The
semicolon, if present, belongs to the body, whose value it discards
(Section~\ref{sec:flow:blocks}).

An \token{else} followed by \token{if} does not start a new construct: the
second conditional is the body of the \token{else}. The two forms below are the
same expression.

\begin{minipage}[t][][t]{0.47\linewidth}
\begin{lstlisting}[style=coloredverbatim, caption=A conditional as the body of \token{else}]
fn foo(c1: bool, c2: bool) {
  if c1 {
    println("c1");
  } else if c2 {
    println("c2");
  } else {
    println("neither");
  }
}
\end{lstlisting}
\end{minipage}\hspace{10pt}%
\begin{minipage}[t][][t]{0.47\linewidth}
\begin{lstlisting}[style=coloredverbatim, caption=A block as the body of \token{else}]
fn foo(c1: bool, c2: bool) {
  if c1 {
    println("c1");
  } else {
    if c2 {
      println("c2");
    } else {
      println("neither");
    }
  }
}
\end{lstlisting}
\end{minipage}

\subsection{Evaluation}

The condition is evaluated first, once. If it is \token{true}, the first body
is evaluated; otherwise the \token{else} body, if any. Exactly one body is
evaluated, or none when the condition is \token{false} and there is no
\token{else}.

The condition must be of type \token{bool}. No other type converts implicitly to
\token{bool}: an integer, a pointer or an option must be compared explicitly.

The condition must not be known at compile time. A condition that the compiler
evaluates to \token{true} makes the \token{else} body dead code, and one it
evaluates to \token{false} makes the first body dead code; both are errors. The
check covers any value the compiler can compute, such as the length of an
array.

\begin{lstlisting}[style=coloredverbatimError, escapechar=@]
let a = [1, 2, 3];
if @\hb{a.len <= 3}@ { // error, always true: the length of an array is known
  println(a[2]);
}

if @\hb{a.len >= 4}@ { // error, always false: the body is dead code
  println(a[0]);
}
\end{lstlisting}

The next listing is a conditional whose condition is a parameter. It lowers to
the L-YIL of Listing~\ref{lst:flow:if_cond_simple}, drawn as a flow graph in
Figure~\ref{fig:flow:if_cond_simple}: the value of the conditional is a variable
(\token{YI\_2}) that each branch assigns before joining the other.

\begin{lstlisting}[style=coloredverbatim]
fn foo(cond: bool)-> i32 {
  if cond {
    89
  } else {
    11
  }
}
\end{lstlisting}

\begin{lstlisting}[style=lyilVerb, caption=Simple condition, label=lst:flow:if_cond_simple]
frame :  _Y4main3fooFbZi32 (let cond(#1) : u8)-> i32 {
    let YI_2(#2) : i32;
#IF cond(#1)
#THEN_GOTO then(#1)
#ELSE_GOTO else(#3);
#LABEL then(#1);
    YI_2(#2) = 89;
#GOTO end(#2);
#LABEL else(#3);
    YI_2(#2) = 11;
#GOTO end(#2);
#LABEL end(#2);
    return YI_2(#2);
}
\end{lstlisting}

\begin{figure}[h]
  \centering
  \begin{adjustbox}{max width=\linewidth}
    \begin{tikzpicture}[node distance=6mm]
      \node[cfterm] (start) {start of foo};
      \node[cfstep, below=of start] (decl) {let YI\_2 : i32;};
      \node[cftest, below=of decl] (test) {cond};
      \node[cfstep] (then) at ($(test) + (-28mm, -16mm)$) {YI\_2 = 89;};
      \node[cfstep] (else) at ($(test) + (28mm, -16mm)$) {YI\_2 = 11;};
      \node[cfstep] (ret) at ($(test) + (0, -30mm)$) {return YI\_2;};
      \node[cfterm, below=of ret] (end) {end of foo};

      \draw[cfflow] (start) -- (decl);
      \draw[cfflow] (decl) -- (test);
      \draw[cfflow] (test.west) -| node[cflab, above, pos=0.4] {true} (then.north);
      \draw[cfflow] (test.east) -| node[cflab, above, pos=0.4] {false} (else.north);
      \draw[cfflow] (then.south) |- (ret.west);
      \draw[cfflow] (else.south) |- (ret.east);
      \draw[cfflow] (ret) -- (end);
    \end{tikzpicture}
  \end{adjustbox}
  \caption{\label{fig:flow:if_cond_simple} The flow graph of
    Listing~\ref{lst:flow:if_cond_simple}. Each branch assigns the value of the
    conditional to \token{YI\_2}, then both paths meet again.}
\end{figure}


\subsection{Value and type}
\label{sec:flow:cond_value_type}

A conditional without \token{else} has no value, and its type is \token{void};
its body must then be of type \token{void} too. A conditional with an
\token{else} has a value when its branches do, and its type is determined as
follows.

\begin{itemize}
\item Breaking branches (Section~\ref{sec:flow:unreachable}) are ignored.
\item The type of the conditional is the type of the first remaining branch.
\item The value of every other remaining branch is implicitly cast to that type.
  A value that cannot be cast is an error.
\end{itemize}

In an \token{else if} chain, the rule applies recursively: the nested
conditional is the value of the \token{else} branch.

\begin{lstlisting}[style=coloredverbatimError, escapechar=@]
fn foo(cond: bool) {
  let a = @\hb{if}@ cond {
    12
  }; // error, no else branch to give a value when /cond/ is false

  let b = if cond {
    12
  } else {
    @\hb{"str"}@
  }; // error, incompatible i32 and [c8]

  let d = if cond {
    12
  } else {
    11u8
  }; // ok, 11u8 is implicitly cast to i32, the type of the first branch

  println(a, b, d);
}
\end{lstlisting}

\subsubsection{Breaking branches}

A branch that ends with \token{return}, \token{throw}, \token{panic},
\token{break} or \token{continue} is ignored by the rule above, so the
conditional of Listing~\ref{lst:flow:if_cond_early_return} has a value, of type
\token{i32}, although its first branch produces none.
Figure~\ref{fig:flow:if_cond_early_exit} shows why: on the path where the
condition holds, control leaves the function before \token{a} is needed.

\begin{figure}[h]
  \centering
  \begin{adjustbox}{max width=\linewidth}
    \begin{tikzpicture}[node distance=6mm]
      \node[cfterm] (start) {start of foo};
      \node[cfstep, below=of start] (decl) {let a : i32;};
      \node[cftest, below=of decl] (test) {cond};
      \node[cfstep] (then) at ($(test) + (-28mm, -16mm)$) {return 89;};
      \node[cfstep] (else) at ($(test) + (28mm, -16mm)$) {a = 10;};
      \node[cfstep, below=of else] (ret) {return a + 9;};
      \node[cfterm] (end) at ($(test) + (0, -44mm)$) {end of foo};

      \draw[cfflow] (start) -- (decl);
      \draw[cfflow] (decl) -- (test);
      \draw[cfflow] (test.west) -| node[cflab, above, pos=0.4] {true} (then.north);
      \draw[cfflow] (test.east) -| node[cflab, above, pos=0.4] {false} (else.north);
      \draw[cfflow] (then.south) |- (end.west);
      \draw[cfflow] (else) -- (ret);
      \draw[cfflow] (ret.south) |- (end.east);
    \end{tikzpicture}
  \end{adjustbox}
  \caption{\label{fig:flow:if_cond_early_exit} The flow graph of
    Listing~\ref{lst:flow:if_cond_early_return}. The path where \token{cond}
    holds leaves the function without reaching \token{a}, so only the other
    branch has to give it a value.}
\end{figure}


\begin{lstlisting}[style=coloredverbatim, caption=Early return in \token{if} condition, label=lst:flow:if_cond_early_return]
fn foo(cond: bool)-> i32 {
  let a = if cond {
    return 89; // exits the function
  } else {
    10
  }; // ok, all non-returning branches have a value

  a + 9
}
\end{lstlisting}

\subsubsection{Mutability of the value}

When the branches have compatible types that differ in mutability, the value of
the conditional takes the most restricted mutability among them. In the next
example the first branch is of type \token{mut [mut i32]} and the second of type
\token{mut [i32]}, so the conditional is of type \token{mut [i32]}, which a
\token{dmut} variable cannot hold.

\begin{lstlisting}[style=coloredverbatimError, escapechar=@]
fn foo(cond: bool) {
  let a = copy [1, 2, 3];
  let dmut @\hb{b}@ = if cond { // error, mutability does not fit in /b/
    copy [2, 3, 4] // mut [mut i32]
  } else {
    a // mut [i32]
  };
  println(b);
}
\end{lstlisting}

A branch that moves borrowed mutable data must say so with \token{alias}, as any
mutable data movement does (Chapter~\ref{chap:memory}). The contract is then
already accepted inside the branch, and is not repeated in front of the
\token{if}.

\begin{lstlisting}[style=coloredverbatim]
fn foo(cond: bool) {
  let dmut a = copy [1, 2, 3];
  let dmut b = if cond {
    alias a // alias is mandatory here
  } else {
    copy [2, 3, 4]
  };
  println(b);
}
\end{lstlisting}

\subsection{Conditional pattern}
\label{sec:flow:cond_pattern}

\token{if let pattern = value} evaluates \token{value}, then tests whether it
fits \token{pattern} (Section~\ref{sec:flow:pattern_matching}). If it does, the
variables the pattern declares are initialized, and the first body is evaluated
with them in scope. If it does not, the \token{else} body is evaluated, and the
pattern's variables do not exist in it.

\begin{figure}[h]
  \centering
  \begin{adjustbox}{max width=\linewidth}
    \begin{tikzpicture}[node distance=6mm]
      \node[cfterm] (start) {start of bar};
      \node[cfstep, below=of start] (call) {YI\_1 = foo();};
      \node[cftest, below=of call] (test) {YI\_1.hasValue};
      \node[cfstep] (then) at ($(test) + (-32mm, -18mm)$) {let x = YI\_1.value;\\println(..., x);};
      \node[cfterm] (end) at ($(test) + (0, -34mm)$) {end of bar};

      \draw[cfflow] (start) -- (call);
      \draw[cfflow] (call) -- (test);
      \draw[cfflow] (test.west) -| node[cflab, above, pos=0.4] {true} (then.north);
      \draw[cfflow] (test.south) -- node[cflab, right] {false} (end.north);
      \draw[cfflow] (then.south) |- (end.west);
    \end{tikzpicture}
  \end{adjustbox}
  \caption{\label{fig:flow:if_cond_let_Ok} The flow graph of
    Listing~\ref{lst:flow:if_cond_let_Ok}. \token{x} is declared only on the
    path where the option holds a value.}
\end{figure}


\begin{lstlisting}[style=coloredverbatim, label=lst:flow:if_cond_let_Ok, caption=Example of conditional pattern]
fn foo()-> i32?;

fn bar() {
  if let Ok(x: i32) = foo() {
    println("Foo returned the value : ", x);
  }
}
\end{lstlisting}

Figure~\ref{fig:flow:if_cond_let_Ok} shows the lowering of
Listing~\ref{lst:flow:if_cond_let_Ok}: \token{x} is created only on the path
where the option holds a value.

The pattern must be refutable. A pattern that always fits is a test known at
compile time (Section~\ref{sec:flow:cte}), and is written as a plain \token{let}
(Section~\ref{sec:memory:pattern_vdecl}). A pattern that never fits is rejected
for the same reason.

\begin{lstlisting}[style=coloredverbatimError, escapechar=@]
let a = (1, 2, 3);

if let @\hb{(x, y, z)}@ = a { // error, the pattern always fits
  println(x, ' ', y, ' ', z);
}
\end{lstlisting}

The value and type rules of a conditional pattern are those of
Section~\ref{sec:flow:cond_value_type}.

\subsubsection{Guard}

A \token{guard}, introduced by \token{\&\&} after the value, is a second
condition, of type \token{bool}, evaluated only if the pattern fits, and with
the pattern's variables in scope. The first body is evaluated if both hold, the
\token{else} body otherwise. A guard is the only way to compare a variable of the
pattern with an existing variable, because an identifier in a pattern always
declares a new variable (Section~\ref{sec:flow:pattern_matching}).

\begin{lstlisting}[style=coloredverbatim, caption=Using a pattern guard]
fn foo()-> i32?;

fn bar(x: i32) {
  if let Ok(y) = foo() && y == x {
    println("foo() returned ", y, ", which is equal to ", x);
  } else {
    println("Either none, or y != x");
  }
}
\end{lstlisting}

Without the guard, the equality test is a second conditional inside the first
body, and the case where it fails needs an \token{else} of its own, as in
Listing~\ref{lst:flow:using_second_cond_no_guard}.

\begin{lstlisting}[style=coloredverbatim, caption=Using a second condition instead of pattern guard, label=lst:flow:using_second_cond_no_guard]
fn foo()-> i32?;

fn bar(x: i32) {
  if let Ok(y) = foo() {
    if y == x {
      println("foo() returned ", y, ", which is equal to ", x);
    } else {
      println("y != x");
    }
  } else {
    println("foo() returned none");
  }
}
\end{lstlisting}

Because \token{\&\&} introduces the guard, neither the value nor the guard can
contain an unparenthesized operator of equal or lower precedence: \token{\&\&},
\token{||}, or an assignment. \token{if let Ok(x) = o \&\& x > 1 \&\& x < 5} is a
syntax error, and is written \token{if let Ok(x) = o \&\& (x > 1 \&\& x < 5)}.

\subsection{Compile-time conditional}
\label{sec:flow:cte_if}

\token{cte if} evaluates its condition at compile time, and keeps only the
selected body. The condition must be known at compile time, and the body not
selected is not validated at all: it may refer to symbols that do not exist or
contain type errors, as long as it is syntactically correct. \token{cte} applies
to the whole \token{else if} chain. It does not apply to \token{if let}.

\begin{lstlisting}[style=coloredverbatim, caption=A compile-time conditional]
fn main() {
  cte if i32::size == 4us {
    println("i32 is four bytes");
  } else {
    let x: i32 = "not validated"; // never checked
    println(x);
  }
}
\end{lstlisting}

A \token{cte if} used as a value has the value and type of the selected body
alone; the other body plays no part in the type.

\subsection{Diagnostics}

\begin{center}
  \begin{adjustbox}{max width=\linewidth}
    \begin{tabular}{|l|l|}
      \hline
      Code & Raised when\\[0pt]
      \hline
      \hline
      E4095 & the condition is not of type \token{bool}, or two branches have incompatible types\\[0pt]
      E4268 & the condition, or the pattern of an \token{if let}, is known to hold\\[0pt]
      E4266 & the condition is known not to hold\\[0pt]
      E4006 & the array pattern of an \token{if let} has a length the value never has\\[0pt]
      E4081 & a conditional without \token{else} has a body that produces a value\\[0pt]
      E4045 & the value does not fit the mutability of its destination\\[0pt]
      E4256 & the condition of a \token{cte if} is not known at compile time\\[0pt]
      \hline
    \end{tabular}
  \end{adjustbox}
\end{center}
